\documentclass[10pt,twocolumn]{article}
\usepackage[T1]{fontenc}
\usepackage[utf8]{inputenc}
\usepackage{lmodern}
\usepackage[margin=1.8cm,top=2.4cm,bottom=2cm,columnsep=0.6cm]{geometry}
\usepackage{fancyhdr}
\usepackage{titlesec}
\usepackage{booktabs}
\usepackage{amsmath,amssymb,amsfonts}
\usepackage{microtype}
\usepackage{xcolor}
\usepackage{mdframed}
\usepackage{parskip}
\usepackage{enumitem}
\usepackage{hyperref}

\definecolor{rulered}{RGB}{180,30,30}
\definecolor{boxgray}{RGB}{245,245,245}
\definecolor{keywordgray}{RGB}{100,100,100}

\pagestyle{fancy}
\fancyhf{}
\fancyhead[L]{\textsc{\small Claude Mag}}
\fancyhead[C]{\small\textit{Machine Learning Research Digest}}
\fancyhead[R]{\small 2026-08-19}
\fancyfoot[C]{\small\thepage}
\renewcommand{\headrulewidth}{0.8pt}

\titleformat{\section}{\normalsize\bfseries\color{rulered}}{}{0pt}{}%
  [\vspace{-4pt}\color{rulered}\rule{\columnwidth}{0.4pt}]
\titlespacing{\section}{0pt}{8pt}{4pt}

\hypersetup{colorlinks=true,urlcolor=rulered,linkcolor=black}

\begin{document}
\setlength{\parindent}{0pt}
\setlength{\parskip}{4pt}

{\small\color{keywordgray}\textbf{Keywords:}
  vision-language-action models \textbullet{}
  credit assignment \textbullet{}
  robotic manipulation \textbullet{}
  GRPO}\par\vspace{4pt}

{\LARGE\bfseries\raggedright
  Temporal GRPO: Beyond Trajectory-Level Credit
  in Vision-Language-Action Reinforcement Learning}\par\vspace{4pt}

{\large\itshape\raggedright
  Decomposing Trajectory Rewards Into Step-Level Signals Improves
  Sample Efficiency and Task Success in Robotic Manipulation
  with VLA Policies}\par\vspace{6pt}

{\small\textbf{By} Yao Zhou, Hang Gao, Fengge Wu,
  Changwen Zheng, Wenwen Qiang}\par\vspace{2pt}

{\small\color{keywordgray}arXiv:2608.13026 \textbullet{} Preprint, August 2026 \textbullet{}
  Robotics (cs.RO)}
\par\vspace{2pt}\noindent{\color{rulered}\rule{\columnwidth}{0.6pt}}\vspace{6pt}

Reinforcement learning for Vision-Language-Action (VLA) models assigns a single
trajectory-level reward to the entire action sequence, treating every step as
equally responsible for success or failure. This coarse credit assignment
conflates critical actions (precise contact) with inconsequential ones (stable
approach), producing an uninformative gradient signal that slows learning.
Temporal GRPO extends Group Relative Policy Optimization with a temporal credit
decomposition mechanism, distributing trajectory rewards to individual steps
based on their temporal context, and demonstrates improved sample efficiency and
task success on robotic manipulation benchmarks.

\begin{mdframed}[backgroundcolor=boxgray,linecolor=rulered,linewidth=1pt,
    innerleftmargin=6pt,innerrightmargin=6pt,
    innertopmargin=6pt,innerbottommargin=6pt]
\textbf{\small AT A GLANCE}\par\vspace{4pt}
\begin{description}[leftmargin=0pt,labelsep=4pt,itemsep=2pt,parsep=0pt,
    font=\small\bfseries]
  \item[Problem:] Trajectory-level RL credit assignment for VLA models fails to
    distinguish the contribution of individual actions to task outcomes.
  \item[Why it matters:] VLA manipulation tasks have long action horizons where
    a few critical steps determine success; coarse credit wastes gradient budget.
  \item[Method:] Decompose trajectory reward into step-level credits conditioned
    on temporal position and local state transitions within the GRPO framework.
  \item[Result:] Improved sample efficiency and task success rate over standard
    GRPO on multi-step robotic manipulation benchmarks.
  \item[Evidence:] Gains replicated across multiple manipulation tasks; ablations
    confirm temporal conditioning is the key driver.
\end{description}
\end{mdframed}

\section{The Big Picture}

Vision-Language-Action (VLA) models are large multimodal policies that take
images of the scene and a language instruction as input and output a sequence
of robot actions. They have achieved impressive results on diverse manipulation
tasks by combining the knowledge of large pretrained vision-language models
with task-specific RL fine-tuning.

The RL fine-tuning phase typically uses algorithms like GRPO, which generates
groups of $G$ trajectories from the current policy, computes a binary reward
(did the task succeed?), and uses the relative ranking within the group to
compute policy gradient updates. Every action in a successful trajectory
receives a positive update; every action in a failed trajectory receives a
negative update.

This trajectory-level assignment is coarse. A grasping task with 80 action
steps may succeed or fail based on a single step — the precise contact between
the gripper and the object. Temporal GRPO identifies that step and concentrates
gradient signal on it, leaving the uninformative steps with near-zero gradient.

\section{Why Existing Approaches Fall Short}

\textbf{Standard GRPO} assigns trajectory reward uniformly across all steps:
\[
  \nabla \mathcal{L} = \frac{1}{G}\sum_{g=1}^G \hat{A}_g \cdot
  \sum_{t=1}^T \nabla \log \pi_\theta(a_{g,t} | s_{g,t})
\]
where $\hat{A}_g = R_g - \text{baseline}$ is the group-relative advantage.
Every step $t$ in trajectory $g$ receives the same weight $\hat{A}_g$.

\textbf{Dense reward shaping} assigns step-level rewards via a hand-designed
function (e.g., distance to goal, object height). This works for simple tasks
but requires domain expertise and does not generalize across task families.

\textbf{Hindsight reward assignment} assigns credit based on outcome
retrospectively but does not distinguish which actions were causally important
for that outcome.

Temporal GRPO avoids hand-designed rewards and operates within the standard
GRPO framework, requiring only a temporal credit decomposition function that
is learned from the trajectory data.

\section{Core Mechanism}

Temporal GRPO decomposes the trajectory-level advantage $\hat{A}_g$ into
step-level advantages $\hat{a}_{g,t}$ using a temporal credit function $c(t, T, s_{g,t})$
that depends on:
\begin{itemize}[itemsep=2pt,parsep=0pt]
  \item \textbf{Temporal position $t/T$:} Early steps (approach) and late
    steps (contact, post-grasp) receive different baseline credit weights.
  \item \textbf{State transition:} The change in scene state $\Delta s_{g,t} = s_{g,t+1} - s_{g,t}$
    identifies steps where significant state changes occurred (contact, placement).
    High-change steps receive higher credit.
\end{itemize}

The group-relative baseline is computed at the step level: for step $t$ across
all $G$ trajectories, the baseline is the mean step-level credit over the group:
\[
  \hat{a}_{g,t} = c(t, T, s_{g,t}) \cdot \hat{A}_g - \bar{c}_t
\]
where $\bar{c}_t = \frac{1}{G}\sum_{g'=1}^G c(t, T, s_{g',t}) \cdot \hat{A}_{g'}$.

\section{Technical Formulation}

The Temporal GRPO policy gradient objective is:
\[
  \mathcal{L}_{\text{TG}}(\theta) = -\frac{1}{G}\sum_{g=1}^G \sum_{t=1}^T
  \hat{a}_{g,t} \cdot \log \pi_\theta(a_{g,t} \mid s_{g,t})
\]
with credit function:
\[
  c(t, T, s_{g,t}) = w_{\text{pos}}(t/T) \cdot w_{\text{trans}}(\|\Delta s_{g,t}\|)
\]
where $w_{\text{pos}}$ is a learned position-dependent weight function (modeled
as a shallow MLP) and $w_{\text{trans}}$ is a monotone function of state
transition magnitude.

The position weight $w_{\text{pos}}$ is trained jointly with $\pi_\theta$ via
a meta-learning objective that maximizes the correlation between step-level
advantage and final task success probability:
\[
  \mathcal{L}_{\text{meta}}(w) = -\text{Corr}\left(\hat{a}_{g,t},\,
  P(\text{success} \mid s_{g,1:t})\right)
\]
where $P(\text{success} \mid s_{g,1:t})$ is estimated by a lightweight success
predictor trained on the same trajectory data.

\section{Learning or Inference Procedure}

Training alternates between two updates:
\begin{enumerate}[itemsep=2pt,parsep=0pt]
  \item \textbf{Policy update:} Sample $G$ trajectories from $\pi_\theta$,
    compute $\hat{a}_{g,t}$ using current $w_{\text{pos}}$, update $\theta$
    via $\mathcal{L}_{\text{TG}}$.
  \item \textbf{Credit function update:} Update $w_{\text{pos}}$ via
    $\mathcal{L}_{\text{meta}}$ using the current trajectory data.
\end{enumerate}
The success predictor is updated with every new trajectory batch. At inference,
the VLA policy $\pi_\theta$ is used directly; no credit function is needed.

\section{What the Guarantee Says}

No formal guarantee is provided. The empirical claim is that temporal credit
decomposition consistently improves sample efficiency (fewer trajectories to
reach a target success rate) and asymptotic task success rate over standard GRPO,
across the tested manipulation tasks. The meta-learning correlation objective
provides a principled basis for learning the credit function, but convergence
properties are empirical.

\section{Experimental Findings}

\begin{itemize}[itemsep=2pt,parsep=0pt]
  \item \textbf{Sample efficiency:} Temporal GRPO reaches 80\% task success
    in 40--55\% fewer environment interactions than standard GRPO across
    grasping, placing, and stacking tasks.
  \item \textbf{Asymptotic success:} Final success rates improve by 5--12
    percentage points on complex long-horizon tasks (4+ subtasks).
  \item \textbf{Long-horizon advantage:} Improvement scales with task horizon.
    Short tasks (1--2 steps) show minimal gain; tasks with 6+ action phases
    show the largest improvements.
  \item \textbf{Contact step analysis:} The credit function concentrates
    $>$60\% of the total trajectory credit on the 2--5 contact steps in a
    typical 80-step grasp trajectory, confirming that temporal credit is
    identifying the causally important steps.
\end{itemize}

\section{Ablations and Interpretation}

\textbf{Position weighting only} (removing state-transition weighting) captures
60--70\% of the full benefit, confirming that temporal position is the dominant
signal. The state-transition term adds the remainder.

\textbf{Fixed vs.\ learned position weight:} A hand-designed position weight
(e.g., linear ramp emphasizing later steps) achieves 80\% of the benefit of the
learned weight on structured manipulation tasks, suggesting the timing pattern is
partially predictable from task structure.

\textbf{Group size $G$:} Temporal GRPO achieves the same accuracy as standard
GRPO with $G=8$ using $G=4$, halving the rollout cost for the same gradient quality.

\section*{Reference}

Yao Zhou, Hang Gao, Fengge Wu, Changwen Zheng, Wenwen Qiang.
\textbf{Temporal GRPO: Beyond Trajectory-Level Credit in Vision-Language-Action
Reinforcement Learning}.
arXiv:2608.13026, August 2026.
\url{https://arxiv.org/abs/2608.13026}

\end{document}
